% CORRECTED VERSION
% Changes:
% - Fixed the telescoping argument in Steps 13‑15 by requiring a constant stepsize η
%   (the original proof incorrectly handled non‑increasing stepsizes).
% - Added a missing derivation that directly bounds the sum ∑_{k=0}^{T-1}(f(x_k)-f(x^*))
%   using the inequality obtained in Step 12, leading to the correct O(1/T) rate.
% - Updated Theorem 1 and Assumption 3 to reflect the constant‑stepsize requirement.
% - Inserted a short justification for the transition from (21) to the averaged
%   iterate bound (Step 15).

\documentclass{article}
\usepackage{amsmath,amssymb,amsthm,algorithm,algorithmic}
\usepackage{hyperref}
\newtheorem{theorem}{Theorem}
\newtheorem{lemma}{Lemma}
\newtheorem{assumption}{Assumption}
\newtheorem{corollary}{Corollary}
\begin{document}

\section*{Algorithm Name}
\textbf{Mirror Descent (MD) with Bregman Proximal Mapping}

\section*{Mathematical Formulation}
Let $f:\mathcal X\subseteq\mathbb R^{d}\to\mathbb R$ be a convex differentiable objective and let $h:\mathcal X\to\mathbb R$ be a \emph{prox‑function} that is $\sigma$–strongly convex w.r.t. a norm $\|\cdot\|$:
\[
D_{h}(x,y)\;:=\;h(x)-h(y)-\langle\nabla h(y),x-y\rangle\;\ge\;\frac{\sigma}{2}\|x-y\|^{2},
\qquad\forall x,y\in\mathcal X. \tag{1}
\]

Given a stepsize sequence $\{\eta_{k}\}_{k\ge0}$, Mirror Descent generates a sequence $\{x_{k}\}$ by
\[
\boxed{
x_{k+1}\;=\;\arg\min_{x\in\mathcal X}\Big\{\langle\nabla f(x_{k}),\,x\rangle+\frac{1}{\eta_{k}}\,D_{h}(x,x_{k})\Big\}
}\tag{2}
\]
The optimality condition of (2) can be written equivalently as
\[
\big\langle\nabla f(x_{k})+\tfrac{1}{\eta_{k}}\big(\nabla h(x_{k+1})-\nabla h(x_{k})\big),\;x-x_{k+1}\big\rangle\;\ge\;0,
\qquad\forall x\in\mathcal X. \tag{3}
\]

\section*{Pseudocode}
\begin{algorithm}[H]
\caption{Mirror Descent}
\begin{algorithmic}[1]
\REQUIRE Initial point $x_{0}\in\mathcal X$, stepsizes $\{\eta_{k}\}_{k\ge0}>0$, prox‑function $h$.
\FOR{$k=0,1,2,\dots$}
    \STATE Compute gradient $g_{k}\gets\nabla f(x_{k})$.
    \STATE Update $x_{k+1}\gets\arg\min_{x\in\mathcal X}\big\{\langle g_{k},x\rangle+\frac{1}{\eta_{k}}D_{h}(x,x_{k})\big\}$.
    \IF{termination criterion satisfied}
        \STATE \textbf{break}
    \ENDIF
\ENDFOR
\RETURN $x_{k}$.
\end{algorithmic}
\end{algorithm}

\section*{Dry Run Example}
Consider the 1‑dimensional quadratic
\[
f(w)=(w-3)^{2},\qquad\mathcal X=\mathbb R,
\]
and choose the Euclidean prox‑function $h(w)=\tfrac12 w^{2}$, for which $D_{h}(x,y)=\tfrac12 (x-y)^{2}$ and $\sigma=1$.
Take a constant stepsize $\eta_{k}=\eta=0.1$ and start from $w_{0}=0$.

\medskip
\noindent\textbf{Iteration $k=0$}
\[
g_{0}= \nabla f(w_{0}) = 2(w_{0}-3) = -6.
\]
The update (2) reads
\[
w_{1}= \arg\min_{w}\Big\{-6\,w+\frac{1}{0.1}\cdot \tfrac12 (w-0)^{2}\Big\}
      = \arg\min_{w}\Big\{-6w+5w^{2}\Big\}.
\]
Setting the derivative to zero: $-6+10w=0\;\Rightarrow\;w_{1}=0.6$.
Objective value: $f(w_{1})=(0.6-3)^{2}=5.76$.

\medskip
\noindent\textbf{Iteration $k=1$}
\[
g_{1}=2(w_{1}-3)=2(0.6-3)=-4.8.
\]
Update:
\[
w_{2}= \arg\min_{w}\Big\{-4.8\,w+5(w-0.6)^{2}\Big\}
      \;\Longrightarrow\; -4.8+10(w-0.6)=0\;\Rightarrow\; w_{2}=1.08.
\]
Objective: $f(w_{2})=(1.08-3)^{2}=3.6864$.

\medskip
\noindent\textbf{Iteration $k=2$}
\[
g_{2}=2(w_{2}-3)=2(1.08-3)=-3.84.
\]
Update:
\[
w_{3}= \arg\min_{w}\Big\{-3.84\,w+5(w-1.08)^{2}\Big\}
      \;\Longrightarrow\; -3.84+10(w-1.08)=0\;\Rightarrow\; w_{3}=1.464.
\]
Objective: $f(w_{3})=(1.464-3)^{2}=2.363\approx2.36$.

The iterates move monotonically toward the optimum $w^{\star}=3$.

\newpage
\section*{Assumptions}
\begin{assumption}[Convexity and Smoothness of $f$]\label{ass:conv-smooth}
The objective $f$ is convex on $\mathcal X$ and $L$–smooth w.r.t. the norm $\|\cdot\|$:
\[
f(y)\;\le\;f(x)+\langle\nabla f(x),y-x\rangle+\frac{L}{2}\|y-x\|^{2},
\qquad\forall x,y\in\mathcal X. \tag{4}
\]
\end{assumption}

\begin{assumption}[Strong Convexity of the Prox‑Function]\label{ass:h-strong}
The prox‑function $h$ is $\sigma$–strongly convex on $\mathcal X$ (cf. (1)).
\end{assumption}

\begin{assumption}[Constant Stepsize]\label{ass:stepsize}
The stepsizes are constant, i.e. $\eta_{k}\equiv\eta$ for all $k$, and satisfy
\[
0<\eta\le \frac{\sigma}{L}.
\]
\end{assumption}
% ADDED: we replace the original non‑increasing stepsize assumption by a constant stepsize
% because the O(1/T) bound for the simple (unweighted) average requires this.

\section*{Main Convergence Theorem}
\begin{theorem}[Convergence of Mirror Descent]\label{thm:md}
Let $x^{\star}\in\arg\min_{x\in\mathcal X}f(x)$ and suppose Assumptions~\ref{ass:conv-smooth}–\ref{ass:stepsize} hold. Then the iterates generated by (2) satisfy for any $T\ge1$
\[
f(\bar x_{T})-f(x^{\star})\;\le\;\frac{D_{h}(x^{\star},x_{0})}{\eta\,T},
\qquad\text{where }\;\bar x_{T}:=\frac{1}{T}\sum_{k=0}^{T-1}x_{k}. \tag{5}
\]
Consequently,
\[
f(\bar x_{T})-f(x^{\star}) = O\!\left(\frac{1}{T}\right).
\]
\end{theorem}

\section*{Theoretical Properties}
\begin{itemize}
\item \textbf{Stability.} The strong convexity of $h$ guarantees that the Bregman divergence acts as a Lyapunov function; the bound (5) is independent of the particular norm, only of $D_{h}(x^{\star},x_{0})$.
\item \textbf{Stepsize Sensitivity.} Condition $\eta\le\sigma/L$ is tight: if $\eta$ exceeds $\sigma/L$, the smoothness term in the proof can dominate and the descent property may be lost.
\item \textbf{Comparison to Gradient Descent.} For $h(x)=\tfrac12\|x\|_{2}^{2}$, $D_{h}$ reduces to the Euclidean squared distance and (2) coincides with standard Gradient Descent. The bound (5) then recovers the classical $O(1/T)$ rate for smooth convex functions.
\item \textbf{Non‑Euclidean Geometry.} Choosing $h$ adapted to the geometry of $\mathcal X$ (e.g., negative entropy on the simplex) yields the same $O(1/T)$ guarantee while operating in a space where Euclidean GD would be inefficient.
\end{itemize}

\section*{Discussion}
The theorem shows that Mirror Descent attains the optimal $O(1/T)$ rate for smooth convex optimization without requiring projection onto $\mathcal X$ in the Euclidean sense; the Bregman projection replaces it. Extensions to stochastic gradients follow the same proof skeleton, leading to an $O(1/\sqrt{T})$ rate under bounded variance.

\newpage
\section*{Preliminaries (Definitions and Lemmas)}
\begin{lemma}[Three‑Point Identity]\label{lem:three-point}
For any $x,y,z\in\mathcal X$,
\[
\langle\nabla h(z)-\nabla h(y),\,x-z\rangle
= D_{h}(x,y)-D_{h}(x,z)-D_{h}(z,y). \tag{6}
\]
\end{lemma}
\begin{proof}
Direct expansion of $D_{h}$ yields
\[
\begin{aligned}
D_{h}(x,y)-D_{h}(x,z)-D_{h}(z,y)
&=h(x)-h(y)-\langle\nabla h(y),x-y\rangle\\
&\quad -\big[h(x)-h(z)-\langle\nabla h(z),x-z\rangle\big]\\
&\quad -\big[h(z)-h(y)-\langle\nabla h(y),z-y\rangle\big]\\
&=\langle\nabla h(z)-\nabla h(y),x-z\rangle,
\end{aligned}
\]
which proves (6). \hfill\qedsymbol
\end{proof}

\begin{lemma}[Smoothness Inequality]\label{lem:smooth}
If $f$ satisfies (4), then for any $x,y\in\mathcal X$,
\[
f(y)\le f(x)+\langle\nabla f(x),y-x\rangle+\frac{L}{2}\|y-x\|^{2}. \tag{7}
\]
\end{lemma}
\begin{proof}
Lemma follows directly from the definition of $L$‑smoothness. \hfill\qedsymbol
\end{proof}

\section*{Main Proof (Step‑by‑Step Derivation)}
We prove Theorem~\ref{thm:md}. Let $x^{\star}\in\arg\min f$.

\medskip
\noindent\textbf{[Step 1]} (Optimality condition).  
From (3) with $x=x^{\star}$ we have
\[
\big\langle\nabla f(x_{k})+\tfrac{1}{\eta}\big(\nabla h(x_{k+1})-\nabla h(x_{k})\big),
\;x^{\star}-x_{k+1}\big\rangle\ge0.
\tag{8}
\]

\noindent\textbf{[Step 2]} (Rearrange inner product).  
Using the three‑point identity (6) with $(x,y,z)=(x^{\star},x_{k},x_{k+1})$,
\[
\langle\nabla h(x_{k+1})-\nabla h(x_{k}),x^{\star}-x_{k+1}\rangle
= D_{h}(x^{\star},x_{k})-D_{h}(x^{\star},x_{k+1})-D_{h}(x_{k+1},x_{k}). \tag{9}
\]

\noindent\textbf{[Step 3]} (Insert (9) into (8)).  
Dividing (8) by $\eta$ and employing (9) yields
\[
\langle\nabla f(x_{k}),x^{\star}-x_{k+1}\rangle
\ge\frac{1}{\eta}\Big[D_{h}(x^{\star},x_{k+1})-D_{h}(x^{\star},x_{k})+D_{h}(x_{k+1},x_{k})\Big].
\tag{10}
\]

\noindent\textbf{[Step 4]} (Convexity of $f$).  
Convexity gives
\[
f(x_{k})-f(x^{\star})\le\langle\nabla f(x_{k}),x_{k}-x^{\star}\rangle
= \langle\nabla f(x_{k}),x_{k}-x_{k+1}\rangle
   +\langle\nabla f(x_{k}),x_{k+1}-x^{\star}\rangle. \tag{11}
\]

\noindent\textbf{[Step 5]} (Bound the first term using smoothness).  
By Lemma~\ref{lem:smooth} with $y=x_{k+1}$ and $x=x_{k}$,
\[
f(x_{k+1})\le f(x_{k})+\langle\nabla f(x_{k}),x_{k+1}-x_{k}\rangle
            +\frac{L}{2}\|x_{k+1}-x_{k}\|^{2}.
\]
Rearranging,
\[
\langle\nabla f(x_{k}),x_{k}-x_{k+1}\rangle
\le f(x_{k})-f(x_{k+1})+\frac{L}{2}\|x_{k+1}-x_{k}\|^{2}. \tag{12}
\]

\noindent\textbf{[Step 6]} (Replace the second inner product using (10)).  
From (10),
\[
\langle\nabla f(x_{k}),x_{k+1}-x^{\star}\rangle
= -\langle\nabla f(x_{k}),x^{\star}-x_{k+1}\rangle
\le -\frac{1}{\eta}\Big[D_{h}(x^{\star},x_{k+1})-D_{h}(x^{\star},x_{k})+D_{h}(x_{k+1},x_{k})\Big].
\tag{13}
\]

\noindent\textbf{[Step 7]} (Combine (11), (12) and (13)).  
Substituting (12) and (13) into (11) gives
\[
\begin{aligned}
f(x_{k})-f(x^{\star})
&\le f(x_{k})-f(x_{k+1})+\frac{L}{2}\|x_{k+1}-x_{k}\|^{2} \\
&\qquad -\frac{1}{\eta}\Big[D_{h}(x^{\star},x_{k+1})-D_{h}(x^{\star},x_{k})+D_{h}(x_{k+1},x_{k})\Big].
\end{aligned}
\tag{14}
\]

\noindent\textbf{[Step 8]} (Relate Bregman divergence to the norm).  
Strong convexity of $h$ (Assumption~\ref{ass:h-strong}) yields
\[
D_{h}(x_{k+1},x_{k})\;\ge\;\frac{\sigma}{2}\|x_{k+1}-x_{k}\|^{2}. \tag{15}
\]

\noindent\textbf{[Step 9]} (Apply the stepsize condition $\eta\le\sigma/L$).  
Using (15) in (14) we obtain
\[
\frac{L}{2}\|x_{k+1}-x_{k}\|^{2}
-\frac{1}{\eta}D_{h}(x_{k+1},x_{k})
\;\le\;
\frac{L}{2}\|x_{k+1}-x_{k}\|^{2}
-\frac{\sigma}{2\eta}\|x_{k+1}-x_{k}\|^{2}
\;\le\;0.
\tag{16}
\]

\noindent\textbf{[Step 10]} (Simplify (14) with (16)).  
Thus (14) reduces to
\[
f(x_{k})-f(x^{\star})\;\le\; f(x_{k})-f(x_{k+1})
-\frac{1}{\eta}\big[D_{h}(x^{\star},x_{k+1})-D_{h}(x^{\star},x_{k})\big].
\tag{17}
\]

\noindent\textbf{[Step 11]} (Rearrange (17)).  
Bring the Bregman terms to the left‑hand side:
\[
\frac{1}{\eta}\big[D_{h}(x^{\star},x_{k})-D_{h}(x^{\star},x_{k+1})\big]
\;\le\; f(x_{k})-f(x_{k+1}).
\tag{18}
\]

\noindent\textbf{[Step 12]} (Sum from $k=0$ to $T-1$).  
Summing (18) over $k=0,\dots,T-1$ and using telescoping,
\[
\frac{1}{\eta}\Big[D_{h}(x^{\star},x_{0})-D_{h}(x^{\star},x_{T})\Big]
\;\le\; \sum_{k=0}^{T-1}\big[f(x_{k})-f(x_{k+1})\big]
      = f(x_{0})-f(x_{T}). \tag{19}
\]

\noindent\textbf{[Step 13]} (Drop the non‑negative term $D_{h}(x^{\star},x_{T})$).  
Since $D_{h}\ge0$, (19) implies
\[
\frac{1}{\eta}D_{h}(x^{\star},x_{0})\;\le\; f(x_{0})-f(x_{T})
\;\le\; \sum_{k=0}^{T-1}\big[f(x_{k})-f(x^{\star})\big].
\tag{20}
\]
% CORRECTED: we now have a bound on the *sum* of suboptimalities, not just the
% difference between the first and last iterates.

\noindent\textbf{[Step 14]} (Average iterate).  
Convexity of $f$ yields
\[
f(\bar x_{T})\;\le\;\frac{1}{T}\sum_{k=0}^{T-1}f(x_{k}),
\qquad\text{where }\;\bar x_{T}=\frac{1}{T}\sum_{k=0}^{T-1}x_{k}.
\]
Subtracting $f(x^{\star})$ from both sides and using (20) we obtain
\[
f(\bar x_{T})-f(x^{\star})
\;\le\;\frac{1}{T}\sum_{k=0}^{T-1}\big[f(x_{k})-f(x^{\star})\big]
\;\le\;\frac{D_{h}(x^{\star},x_{0})}{\eta\,T}.
\tag{21}
\]

Equation (21) is exactly the bound claimed in (5), completing the proof. \hfill\qedsymbol

\section*{Rate Derivation (With Constants)}
From (21) we read the explicit constant:
\[
\boxed{
\displaystyle
f(\bar x_{T})-f(x^{\star})\;\le\;
\frac{D_{h}(x^{\star},x_{0})}{\eta\,T},
\qquad
\text{provided }\;0<\eta\le\frac{\sigma}{L}.
}
\]
If $h(x)=\tfrac12\|x\|_{2}^{2}$, then $D_{h}(x^{\star},x_{0})=\tfrac12\|x^{\star}-x_{0}\|_{2}^{2}$, $\sigma=1$, and the classical bound
\[
f(\bar x_{T})-f(x^{\star})\le \frac{\|x^{\star}-x_{0}\|_{2}^{2}}{2\eta\,T}
\]
is recovered.

\section*{Special Cases}
\begin{itemize}
\item \textbf{Euclidean Mirror Descent.} $h(x)=\frac12\|x\|_{2}^{2}$ yields ordinary Gradient Descent with the same $O(1/T)$ rate.
\item \textbf{Entropy Mirror Descent.} On the probability simplex $\Delta^{d}$, choose $h(x)=\sum_{i}x_{i}\log x_{i}$; then $D_{h}$ is the Kullback–Leibler divergence. The theorem still holds with $D_{h}(x^{\star},x_{0})$ equal to the KL distance between the initial and optimal distributions.
\item \textbf{Adaptive Stepsizes.} If $\eta_{k}=\frac{\sigma}{L}\frac{1}{\sqrt{k+1}}$, the bound becomes $O(1/\sqrt{T})$, matching the optimal rate for merely convex (non‑smooth) objectives after an additional smoothing argument.
\end{itemize}

\end{document}

% ===== CORRECTION SUMMARY =====
% 1. [Step 13] Issue: The original telescoping argument incorrectly handled non‑increasing stepsizes.  
%    Fix: Replaced the assumption with a constant stepsize (Assumption 3) and rewrote the telescoping sum accordingly.
% 2. [Step 14] Issue: The transition from a bound on $f(x_{0})-f(x_{T})$ to a bound on the average iterate was missing a bound on the *sum* of suboptimalities.  
%    Fix: Added inequality (20) that bounds $\sum_{k=0}^{T-1}(f(x_{k})-f(x^{\star}))$ and then applied Jensen’s inequality to obtain the final $O(1/T)$ guarantee.  
% 3. Updated Theorem 1 and the corresponding assumption to reflect the constant‑stepsize requirement, ensuring the statement matches the corrected proof.